\section{Introduction}
\label{sec:intro}

Shear wall systems are the principal lateral-force-resisting components in many high-rise residential buildings, particularly in seismic regions. The arrangement of shear walls influences stiffness distribution, load paths, and construction cost, and it must simultaneously respect architectural constraints such as functional zoning, circulation, and openings. In preliminary design, these competing requirements are commonly reconciled through iterative manual layout adjustment and repeated coordination between architectural and structural teams. Such a workflow is labor-intensive and tends to restrict systematic exploration of alternative layouts.

Despite extensive design experience and code provisions, routine shear wall layout development remains inefficient for three reasons. First, the mapping from architectural drawings to structurally rational wall candidates relies on implicit, experience-driven rules that are difficult to formalize and reuse. Second, resolving spatial conflicts often requires multiple cross-discipline iterations, creating a slow feedback loop during early design. Third, the quality and consistency of outcomes can vary across designers and projects, while the limited ability to explore the design space can lead to locally satisfactory but globally suboptimal layouts. These limitations motivate automated methods that can interpret architectural constraints and generate feasible wall layouts with controllable design intent.

Recent data-driven studies have treated layout generation as an image-to-image translation task, using convolutional networks, conditional generative adversarial networks, or diffusion models to map rasterized floor plans to structural layouts. While pixel-based representations can capture complex spatial patterns, they introduce inefficiencies that are particularly pronounced for structural elements: thin wall lines occupy a small fraction of the image grid, leading to high redundancy; topological relationships are only implicitly represented; and raster outputs require non-trivial post-processing to recover vectorized wall geometry for downstream analysis.

To address these issues, graph neural networks (GNNs) have been adopted to model structural topology more explicitly. A common formulation builds graphs from geometric primitives, with nodes as wall intersections and edges as wall segments, and predicts whether each segment should be a shear wall or the proportion of shear wall at both ends of the segment. This representation preserves connectivity but remains low-level: individual segments carry limited information about the architectural spaces they bound. As a result, room-level constraints (e.g., functional requirements and openings) are not directly expressed, and related structural layout tasks may require separate representations or additional prediction stages.

A room-level representation is considered more aligned with the physical and architectural constraints governing shear wall placement. In typical residential plans, feasible shear wall candidates are largely restricted to room boundaries, and many practical constraints (openings, circulation, functional usage) are naturally described at the room level rather than at the level of fragmented wall primitives. Modeling rooms as graph nodes and room adjacencies as edges therefore provides an explicit interface between architectural semantics and structural layout generation. In addition, the number of rooms per plan is usually far smaller than the number of primitive wall segments, which reduces graph scale and can improve computational efficiency. The same boundary-based formulation also provides a consistent basis for extending to other boundary-aligned elements such as beams, enabling a more unified structural layout pipeline.

Based on this representation, conditional shear wall layout generation is formulated using a GNN backbone with condition modulation. Neighborhood aggregation is implemented with an attention-based message passing network (e.g., GATv2), while Feature-wise Linear Modulation (FiLM) is used to inject design conditions corresponding to target wall-density categories. A practical challenge is that each floor plan is typically paired with only one engineered layout, limiting condition-specific supervision. To mitigate this, training is organized into a supervised stream that learns geometric fidelity from available layouts and a counterfactual stream that enforces condition compliance under alternative conditions through density-oriented constraints.

The main contributions are as follows:

\begin{itemize}
\item A room-level graph formulation for structural layout generation that encodes architectural spaces and their adjacencies, enabling direct use of room semantics and boundary constraints while reducing the scale of graph inference relative to primitive-based formulations.
\item A conditional GNN architecture with FiLM-based modulation for controllable shear wall layout generation under prescribed wall-density categories, including a decoding design that predicts both wall existence and geometry.
\item A dual-stream training scheme that improves conditional compliance under limited paired data by combining supervised learning with counterfactual, density-oriented regularization.
\end{itemize}

The remainder of this paper is organized as follows. Section 2 reviews related work on deep learning for structural layout generation, GNN-based representations in architecture-engineering-construction, and conditional generation. Section 3 presents the methodology, including the room-level graph construction, conditional model architecture, and training strategy. Section 4 describes the experimental setup and reports quantitative and qualitative results, including ablation analysis and case studies. Section 5 discusses implications and limitations. Section 6 concludes the paper.
